\abstractEN % Do NOT modify this line

The goal of this project is to develop an application that performs operations on a relational database provided by the instructors. In addition to the initial database, it was necessary to implement additional SQL features, namely triggers, a view that supports insert and update operations, and a stored procedure.
The application was implemented in Java, using the Jakarta Persistence API (JPA) for persistence management.
JPA is a standardized API of the Jakarta EE platform that facilitates the mapping between Java objects and relational tables. It allows persistence operations (insertion, update, deletion, query) to be performed through an object-oriented model, abstracting away the direct use of SQL. Furthermore, it provides support for transactions, JPQL queries, and concurrency control through mechanisms such as optimistic locking.
